\chapter{\(3\) hojas temporales, cálculo nilpotente y conexión tritificada de Cartan}
\label{chap:tres-hojas-variacion}

\section{Familia cuadr\'atica de TRIT}

Para \(\tau\in\{+1,0,-1\}\), consideremos el \`algebra real

\[
\mathbb A_\tau=\mathbb R[J_\tau]/(J_\tau^2+\tau),
\qquad
J_\tau^2=-\tau I.
\]

La exponencial se calcula exactamente:

\[
e^{tJ_{+1}}=\cos t+J_{+1}\sin t,
\]

\[
e^{tJ_0}=I+tJ_0,
\]

\[
e^{tJ_{-1}}=\cosh t+J_{-1}\sinh t.
\]

Las tres hojas poseen funciones geom\'etricas diferenciadas:

\begin{center}
\begin{tabular}{ccl}
\toprule
\(\tau\)&\(J_\tau^2\)&r\'egimen\\
\midrule
\(+1\)&\(-I\)&el\'iptico: cierre, retorno y memoria\\
\(0\)&\(0\)&parab\'olico: tangencia, variaci\'on y presente\\
\(-1\)&\(+I\)&hiperb\'olico: apertura, boost y propagaci\'on\\
\bottomrule
\end{tabular}
\end{center}

Sea \(\widehat\tau=P_+-P_-\), con \(P_0=I-P_+-P_-\). Entonces

\[
\widehat\tau^3=\widehat\tau.
\]

Una involuci\'on que intercambia las hojas extrema y conserva la hoja central
satisface

\[
\mathsf J\widehat\tau\mathsf J^{-1}=-\widehat\tau.
\]

El umbral permanece fijo y las orientaciones de cierre y apertura se
intercambian.

\section{La hoja nula como 1.er jet}

En \(\mathbb A_0\), escribamos \(\epsilon=J_0\), con \(\epsilon^2=0\). Para
todo funcional diferenciable \(F\),

\begin{equation}
\boxed{F(x+\epsilon v)=F(x)+\epsilon\,dF_x(v).}
\label{eq:rm-dual-first-jet}
\end{equation}

La igualdad es exacta porque todos los t\'erminos de orden al menos dos
contienen \(\epsilon^2\). En particular, para una acci\'on
\(S[e,\omega,\Psi]\),

\begin{equation}
S[e+\epsilon\delta e,\omega+\epsilon\delta\omega,
  \Psi+\epsilon\delta\Psi]
=S[e,\omega,\Psi]+\epsilon\,\delta S.
\label{eq:rm-action-first-jet}
\end{equation}

El coeficiente nilpotente contiene las ecuaciones de Euler--Lagrange, la
ecuaci\'on de Cartan, el potencial simpl\'ectico y los t\'erminos de borde.
La hoja central es, por tanto, la residencia algebraica del principio
variacional.

\begin{theorem}[Estacionariedad nilpotente]
Para una configuraci\'on \(z\), son equivalentes:
\[
dS_z(\delta z)=0\quad\text{para toda }\delta z,
\]
y
\[
S(z+\epsilon\delta z)=S(z)
\quad\text{en }\mathbb A_0
\quad\text{para toda }\delta z.
\]
\end{theorem}

\begin{proof}
La equivalencia es la igualdad de los coeficientes de \(1\) y \(\epsilon\) en
\eqref{eq:rm-dual-first-jet}.
\end{proof}

El Hamiltoniano global, la acci\'on efectiva de la hoja central y el
Hamiltoniano modular del borde son objetos coordinados:

\[
H_{\rm glob},
\qquad
L_{0,\rm eff},
\qquad
K_{\rm HHM}=-\log\rho_{\partial}.
\]

El primero genera evoluci\'on; el segundo publica la variaci\'on local; el
tercero conserva procedencia informacional.

\section{Reducci\'on de Feshbach--Schur}

Sea

\[
\mathcal H=\mathcal H_0\oplus\mathcal H_{\rm ext},
\qquad
H=
\begin{pmatrix}
H_{00}&V\\ V^*&H_{\rm ext}
\end{pmatrix}.
\]

Para \(z\) en el conjunto resolvente de \(H_{\rm ext}\), la compresi\'on del
resolvente es

\begin{equation}
P_0(z-H)^{-1}P_0
=\left[z-H_{00}-\Sigma_0(z)\right]^{-1},
\qquad
\Sigma_0(z)=V(z-H_{\rm ext})^{-1}V^*.
\label{eq:rm-feshbach-schur}
\end{equation}

El operador \(\Sigma_0\) transporta a la hoja presente la memoria de las hojas
el\'iptica e hiperb\'olica. Su transformada temporal genera un n\'ucleo de
memoria

\[
S_{\rm eff}[\psi_0]
=\int\psi_0^\dagger(i\hbar\partial_t-H_{00})\psi_0\,\mathrm dt
-\iint\psi_0^\dagger(t)K_0(t,t')\psi_0(t')\,\mathrm dt\,\mathrm dt'.
\]

Una expansi\'on regular de baja frecuencia,

\[
\Sigma_0(\omega)
=\Sigma_0^{(0)}+\omega\Sigma_0^{(1)}
+\omega^2\Sigma_0^{(2)}+\cdots,
\]

produce una acci\'on local efectiva; una estructura espectral singular conserva
memoria no local. El presente recibe as\'i la respuesta del estado global sin
introducir un tiempo exterior.

\section{Conexi\'on universal de Cartan}

Sea \(\mathfrak g_\tau\) el \`algebra de Cartan generada por
\(J_{ab}=-J_{ba}\) y \(P_a^{(\tau)}\), con relaciones

\begin{align}
 [J_{ab},J_{cd}]&=
 \eta_{bc}J_{ad}-\eta_{ac}J_{bd}
 -\eta_{bd}J_{ac}+\eta_{ad}J_{bc},\\
 [J_{ab},P_c^{(\tau)}]&=
 \eta_{bc}P_a^{(\tau)}-\eta_{ac}P_b^{(\tau)},\\
 [P_a^{(\tau)},P_b^{(\tau)}]&=-\tau J_{ab}.
 \label{eq:rm-cartan-brackets}
\end{align}

Para \(\tau=0\) se obtiene el \`algebra de Poincar\'e; los signos extremos
producen las dos formas de curvatura constante. Sea \(\omega\) una conexi\'on
de Lorentz, \(e=e^aP_a^{(\tau)}\) una cot\'etrada valorada en el sector de
traslaciones y \(\ell\) una escala de longitud. Definimos la conexi\'on
tritificada

\begin{equation}
\mathcal A_\tau^{\rm Cart}
=\frac12\omega^{ab}J_{ab}+\ell^{-1}e^aP_a^{(\tau)}.
\label{eq:rm-trit-cartan-connection}
\end{equation}

Su curvatura es

\begin{equation}
\boxed{
\mathcal F_\tau^{\rm Cart}
=\frac12\left(F_\omega^{ab}-\tau\ell^{-2}e^a\wedge e^b\right)J_{ab}
+\ell^{-1}T^aP_a^{(\tau)},
\qquad T^a=D_\omega e^a.}
\label{eq:rm-trit-cartan-curvature}
\end{equation}

\begin{proof}
Al expandir
\(\mathrm d\mathcal A_\tau^{\rm Cart}
+\mathcal A_\tau^{\rm Cart}\wedge\mathcal A_\tau^{\rm Cart}\), el primer
corchete de \cref{eq:rm-cartan-brackets} produce \(F_\omega\), el segundo
produce \(D_\omega e\) y el tercero aporta
\(-\tau\ell^{-2}e^a\wedge e^bJ_{ab}/2\).
\end{proof}

En la hoja central,

\[
\mathcal F_0^{\rm Cart}
=\frac12F_\omega^{ab}J_{ab}+\ell^{-1}T^aP_a^{(0)}.
\]

La tors\'ion es el coeficiente nilpotente de la curvatura. Esta igualdad une
la interpretaci\'on variacional de la hoja \(0\) con su interpretaci\'on
geom\'etrica de Cartan.

\section{Corona non\'adica y energ\'ia torsional}

Sobre las aristas de la corona de nueves, sea

\[
(D_{9,m}f)(a)=f(t(a))-U_a f(s(a)),
\]

donde \(U_a\) transporta orientaci\'on y acarreo. Para un peso positivo \(W_m\),

\begin{equation}
H_m=H_{{\rm clk},m}+D_{9,m}^*W_mD_{9,m},
\qquad
\mathcal E_{{\rm tor},m}(f)
=\|W_m^{1/2}D_{9,m}f\|^2.
\label{eq:rm-nonadic-torsion-energy}
\end{equation}

La holonom\'ia unitaria conserva positividad. Los cuadrados de refinamiento

\[
D_{9,m+1}J_m=K_mD_{9,m},
\qquad
K_m^*W_{m+1}K_m=W_m
\]

garantizan la compatibilidad de la energ\'ia a toda profundidad.

El cierre posterior construir\'a la soldadura de pantalla que transforma este
defecto en una respuesta de Cartan. La presente secci\'on fija el dominio
operatorio y la positividad que ese entrelazador debe preservar.

\section{Principio de Ambivalencia: involución bidireccional de las razones \(\pi/\varphi^2\) y \(\varphi/\pi^2\)}

Definimos

\[
\mathcal Q(x,y)=\left(\frac{x}{y^2},\frac{y}{x^2}\right),
\qquad x,y>0.
\]

Para

\[
P=xy,
\qquad
O=x/y,
\]

se verifica

\[
P(\mathcal Q(x,y))=P^{-1},
\qquad
O(\mathcal Q(x,y))=O^3.
\]

Después de dos pasos,

\begin{equation}
\boxed{P(\mathcal Q^2)=P,
\qquad O(\mathcal Q^2)=O^9.}
\label{eq:rm-qmap-nine}
\end{equation}

En coordenadas \(u=\log(xy)\), \(v=\log(x/y)\),

\[
(u,v)\longmapsto(-u,3v)\longmapsto(u,9v).
\]

El producto retorna y la orientaci\'on multiplica por nueve su memoria. Para
\((x,y)=(\pi,\varphi)\), las dos salidas son

\[
\rho_A=\frac{\pi}{\varphi^2},
\qquad
\rho_E=\frac{\varphi}{\pi^2}.
\]

Satisfacen

\[
\rho_E\rho_A^2=\varphi^{-3},
\qquad
M_9=(\rho_E\rho_A^2)^6=\varphi^{-18}.
\]

La relaci\'on vincula lectura areal, exponente electr\'onico, retorno non\'adico
y memoria cuadr\'atica sin identificar entre s\'i sus estructuras algebraicas.

\begin{proofstatusbox}
Las \`algebras cuadr\'aticas, el c\'alculo dual, la reducci\'on de Schur, la
curvatura \eqref{eq:rm-trit-cartan-curvature} y la identidad
\eqref{eq:rm-qmap-nine} son exactas. El l\'imite de la energ\'ia torsional y su
realizaci\'on como acci\'on completa de Einstein--Cartan requieren los
cuadrados de refinamiento y el morfismo f\'isico desarrollados despu\'es.
\end{proofstatusbox}

\begin{falsifierbox}
Refutan el cierre interno: un t\'ermino distinto de
\(-\tau\ell^{-2}e\wedge e\) en la curvatura; un coeficiente nilpotente que no
sea \(D_\omega e\); un residuo en la identidad de Feshbach--Schur; p\'erdida
de positividad de \(D_9^*WD_9\); o fallo de cualquiera de las dos ecuaciones
de \eqref{eq:rm-qmap-nine}.
\end{falsifierbox}
